\section{DERIVATIONS}
\label{app:derivations}

\subsection{The Reconstruction Bound}
\label{app:bound}

Fix a cell $i$ and hold $\vc_i$ and $\rhobar_i$ constant, as \cref{sec:batching} does; every probability in this appendix is conditional on $\rhobar_i$, which we suppress from the notation. Then
\begin{align}
  \log p(\vx_i \mid \vc_i, t_i)
  &= \log \iint p(\vx_i \mid \vz, \vw)\, p(\vw \mid \vc_i, t_i)\, p(\vz)\, d\vz\, d\vw \\
  &\geq \E_{q(\vz)\,q(\vw \mid \vz)}\Big[\log p(\vx_i \mid \vz, \vw) + \log p(\vw \mid \vc_i, t_i) + \log p(\vz) - \log q(\vz) - \log q(\vw \mid \vz)\Big] \label{eq:jensen}\\
  &= \E_{q(\vz)}\Big[\E_{q(\vw \mid \vz)}\big[\log p(\vx_i \mid \vz, \vw)\big] - \KL\big(q(\vw \mid \vz)\,\|\,p(\vw \mid \vc_i, t_i)\big)\Big] - \KL\big(q(\vz)\,\|\,p(\vz)\big),
\end{align}
where \cref{eq:jensen} is Jensen's inequality under the variational family $q(\vz)\,q(\vw \mid \vc_i, t_i, \vz, \vx_i)$. Because $q(\vw \mid \cdot)$ conditions on the sampled $\vz$, the exact bound keeps the $\vw$-divergence inside $\E_{q(\vz)}$. The divergence between two diagonal Gaussians is available in closed form,
\begin{equation}
  \KL\big(\N(\vmu, \vsigma^2)\,\|\,\N(\mathbf{m}, \mathbf{s}^2)\big) = \sum_k \Big[\log\frac{s_k}{\sigma_k} + \frac{\sigma_k^2 + (\mu_k - m_k)^2}{2 s_k^2} - \frac{1}{2}\Big],
  \label{eq:gauss-kl}
\end{equation}
and evaluating it at the single reparameterised draw of $\vz$ used for the reconstruction term is an unbiased one-sample estimate of $\E_{q(\vz)}[\KL(\cdot)]$.

Summed over cells, with each neighbour's $\vz_j, \vw_j$ drawn from its own posterior as in \cref{sec:batching}, $\sum_i L_i$ is a one-sample estimate of the evidence lower bound of the coupled model, $\log p(\{\vx_i\} \mid \{\vc_i, t_i\})$, under the fully factorised family $\prod_i q(\vz_i)\, q(\vw_i \mid \vz_i)$. The per-cell bounds are not independent terms, since $\rhobar_i$ is a function of the neighbours' latents. The stop-gradient removes the terms $i \ne j$ from the gradient with respect to $q_j$, so training ascends a partial gradient of this bound: a fixed-point iteration in which each step treats the current $\rhobar$ as data.

\subsection{The Intrinsic Path and the Factor \texorpdfstring{$(1+\omega)$}{(1+omega)}}
\label{app:pathb}

A second bound on the same evidence holds by restricting the variational family to $q(\vw) := p(\vw \mid \vc_i, t_i)$, which zeroes the $\vw$-divergence and leaves
\begin{equation}
  \log p(\vx_i \mid \vc_i, t_i) \geq \E_{q(\vz)}\E_{p(\vw \mid \vc_i, t_i)}\big[\log p(\vx_i \mid \vz, \vw)\big] - \KL\big(q(\vz)\,\|\,p(\vz)\big) =: L_i^{z}.
\end{equation}
Call the bound of \cref{app:bound} $L_i$. Each carries its own $-\KL(q(\vz_i)\|p(\vz_i))$, so $L_i + \omega L_i^{z}$ contains $-(1+\omega)\KL(q(\vz_i)\|p(\vz_i))$. Omitting the second copy breaks the correspondence with $L_i + \omega L_i^{z}$ even at $\omega = 1$.

The construction mirrors the two-bound objective of \citet{slavutsky2025}, where the sum of the shared-only and joint bounds corresponds to minimising $\KL(q(\vz) \,\|\, p(\vz \mid \vx)) + \KL(q(\vz)q(\vw) \,\|\, p(\vz,\vw \mid \vx, u))$ with $u$ their condition label, a sum of two non-negative divergences. Here, unlike in \citet{slavutsky2025}, whose shared-only term decodes from $\vz$ alone and so bounds a different model, both terms bound the same evidence,
\begin{equation*}
  L_i + \omega L_i^{z} \le (1+\omega)\log p(\vx_i \mid \vc_i, t_i),
\end{equation*}
with gap $\KL\big(q(\vz)q(\vw \mid \vz)\,\|\,p(\vz,\vw \mid \vx_i,\vc_i,t_i)\big) + \omega\,\KL\big(q(\vz)p(\vw \mid \vc_i,t_i)\,\|\,p(\vz,\vw \mid \vx_i,\vc_i,t_i)\big)$. One $q(\vz)$ serves two variational families, so its optimum is a compromise and not the best approximation of the posterior under either, which is one of the reasons uncertainty is reported from the leakage sweep rather than from $q$. The invariance term sits outside the bound for a separate reason (\cref{app:posterior-reg}).

Why the intrinsic path is load-bearing: in \citet{slavutsky2025} the shared code is kept informative by a reconstruction from $\vz$ alone and by a prior on $\vw$ centred on the class-wise means of $\vz$. Our intrinsic path is the analogue of the first. We have no analogue of the second: our prior on $\vw$ never references $\vz$, and the prior on $\vz$ is a standard normal, so nothing else in the objective rewards $\vz$ for being informative. The intrinsic path supplies that pressure directly, since under it $\vz$ must explain the cell given only the expected response. It runs through the same leakage mixture as the first term, so it never asserts $\kappa = 0$ and contradicts it.

\subsection{The Invariance Term as Posterior Regularisation}
\label{app:posterior-reg}

Neither $\Pen$ nor $\Adv_i$ is part of any bound. Both are posterior regularisation in the spirit of \citet{ganchev2010}: a constraint on the variational family, $\I(\vz; [\vy, \vPhi] \mid t) = 0$, written as a penalty on the variational objective. This is the same status \citet{slavutsky2025} give their adversarial term. It is a statement about which $q$ are admissible, not a modification of the generative model, which does not model $\vy$ or $\vPhi$ at all; the constraint states what the decomposition is meant to satisfy, and it acts on the aggregate of the encoder's outputs over cells rather than on each cell's posterior, so it is posterior regularisation in a looser sense than the per-instance expectation constraints of \citet{ganchev2010}.

\subsection{The Gaussian Conditional Mutual Information}
\label{app:gaussian-mi}

For jointly Gaussian $(\vz, \vv)$ with covariance $\Sig_{[\vz,\vv]}$ and marginal covariances $\Sig_{\vz}, \Sig_{\vv}$, the mutual information is $\I(\vz;\vv) = \tfrac12[\logdet\Sig_{\vz} + \logdet\Sig_{\vv} - \logdet\Sig_{[\vz,\vv]}]$ \citep{cover2006}. \cref{eq:penalty} applies this per type and averages with the type frequencies, which is the Gaussian estimate of $\I(\vz;\vv \mid t)$ under a type-wise Gaussian assumption. Two properties used in \cref{sec:invariance} follow directly. First, the value is invariant to a per-dimension rescaling $\vu \mapsto \mathbf{D}\vu$ with $\mathbf{D}$ diagonal: each log-determinant gains $\log\det$ of the corresponding block of $\mathbf{D}^2$, and the blocks of $\mathbf{D}$ for $\vz$ and $\vv$ together make up $\mathbf{D}$, so the three terms cancel exactly. The penalty may therefore be computed on correlation matrices. Second, on the simplex $\sum_k y_k = 1$ the covariance of the full composition vector is singular and the log-determinant is $-\infty$; dropping one column removes the constraint without changing the information, since the dropped column is a deterministic function of the others.

\subsection{Feedback Amplification Through the Leakage Path (Heuristic)}
\label{app:amplification}

The following is a first-order heuristic, not a theorem. Write the decoder output of cell $j$ as $\vrho_j$ and suppose gradients were allowed to flow from cell $i$'s reconstruction through $\rhobar_i = \sum_j \beta_{ij}\vrho_j$ into the parameters that produce $\vrho_j$. Linearising around a fixed point of the coupled system, a perturbation $\delta$ of a cell's clean composition induces a perturbation $\kappa\,\beta\,\delta$ of its neighbours' fitted compositions, which feeds back through the same path; the total response is $\delta\sum_{n \ge 0}(\kappa \mathbf{P}_\beta)^n$ for a row-stochastic $\mathbf{P}_\beta$ built from $\beta$, whose spectral radius is at most one, so the geometric series is bounded by $1/(1-\kappa)$. The stop-gradient truncates the series at $n = 0$. The argument concerns perturbations of the fitted compositions; the gradient fixed point through shared decoder parameters is more involved, and we have not fitted the open-path model; \citet{ergen2025} close the same path for training stability. \todo{either complete the gradient-level argument or keep this as a heuristic}
